\Needspace{6\baselineskip}
\begin{center}\hypertarget{native-261-article188}{}\pdfbookmark[2]{188.}{article-188}188.\end{center}
Methodus praecedens iis potissimum casibus adaptata est, ubi praecisio
summa desideratur: saepissime autem occurrunt casus, ubi sine haesitatione
paullulum ab illa remitti potest, si hoc modo calculi prolixitatem considerabiliter
contrahere licet, praesertim quando observationes magnum temporis
intervallum nondum includunt, adeoque de orbitae determinatione ut sic dicam
definitiva nondum cogitatur. In talibus casibus methodus sequens lucro notabili
in usum vocari poterit.

Eligantur e tota observationum copia duo loca completa $L$ et $L'$, computenturque
pro temporibus respondentibus ex elementis approximatis corporis
coelestis distantiae a terra. Formentur dein respectu harum distantiarum
tres hypotheses, retentis in prima valoribus computatis, mutataque in hypothesi
secunda distantia prima, secundaque in hypothesi tertia; utraque mutatio
pro ratione incertitudinis, quae in illis distantiis remanere praesumitur, ad
lubitum accipi poterit. Secundum has tres hypotheses, quas in schemate sequente
exhibemus:
\begin{center}
\begingroup
\setlength{\tabcolsep}{4pt}
\renewcommand{\arraystretch}{1.3}
\begin{tabular}{|l|c|c|c|}
\hline
 & Hyp. I & Hyp. II & Hyp. III\\
\hline
Distantia\footnotemark\ loco primo respondens & $D$ & $D+\delta$ & $D$\\
Distantia loco secundo respondens & $D'$ & $D'$ & $D'+\delta'$\\
\hline
\end{tabular}
\endgroup
\end{center}
\footnotetext{Adhuc commodius erit, loco distantiarum ipsarum logarithmis distantiarum curtatarum uti.}
\noindent\SourceLeaf{262}computentur e duobus locis $L,L'$ per methodos in Libro primo explicatas
tria elementorum systemata, ac dein ex his singulis loca geocentrica corporis
coelestis temporibus omnium reliquarum observationum respondentia. Sint
haec (singulis longitudinibus et latitudinibus, vel ascensionibus rectis et declinationibus
seorsim denotatis)
\begin{center}
\begin{tabular}{@{}l@{\quad}l@{\quad}l@{\quad}l@{\quad}l@{}}
\makebox[12.8em][l]{in systemate primo\hfill .\,.\,.\,.\,.\,.} & $M$, & $M'$, & $M''$ & etc.\\[.3em]
\makebox[12.8em][l]{in systemate secundo\hfill .\,.\,.\,.\,.} & $M+a$, & $M'+a'$, & $M''+a''$ & etc.\\[.3em]
\makebox[12.8em][l]{in systemate tertio\hfill .\,.\,.\,.\,.\,.\,.} & $M+\beta$, & $M'+\beta'$, & $M''+\beta''$ & etc.
\end{tabular}
\end{center}
\noindent Sint porro resp.
\begin{center}
\begin{tabular}{@{}l@{\quad}l@{\quad}l@{\quad}l@{\quad}l@{}}
\makebox[12.8em][l]{loca observata\hfill .\,.\,.\,.\,.\,.\,.\,.\,.\,.\,.} & $N$, & $N'$, & $N''$ & etc.
\end{tabular}
\end{center}

Iam quatenus mutationibus parvis distantiarum $D,D'$ respondent mutationes
proportionales singulorum elementorum, nec non locorum geocentricorum
ex his computatorum, supponere licebit, loca geocentrica e quarto elementorum
systemate computata, quod distantiis a terra $D+x\delta,D'+y\delta'$ superstructum
sit, resp. fore $M+ax+\beta y,M'+a'x+\beta'y,M''+a''x+\beta''y$ etc. Hinc
dein, secundum disquisitiones praecedentes, quantitates $x,y$ ita determinabuntur,
ut illae quantitates cum $N,N',N''$ etc. resp. quam optime consentiant (ratione
praecisionis relativae observationum habita). Systema elementorum correctum
ipsum vel perinde ex $L,L'$ et distantiis $D+x\delta,D'+y\delta'$, vel secundum
regulas notas e tribus elementorum systematibus primis per simplicem
interpolationem derivari poterit.
